
@@M0@@ \uline{de} @@M1@@, \uline{et une suite de sections} @@M2@@ \uline{de} @@M3@@ \uline{telles que le complexe de Koszul associé à} @@M4@@ \uline{soit une résolution de} @@M5@@, \uline{le complexe de Koszul associé à} @@M6@@ \uline{une résolution de} @@M7@@, \uline{et le complexe de Koszul associé à} @@M8@@ \uline{une résolution de} @@M9@@ (\uline{conditions satisfaites en particulier dans le cas noethérien si la suite} @@M10@@ \uline{est régulière et si} @@M11@@, @@M12@@). \uline{Soient} @@M13@@ \uline{les idéaux de} @@M14@@. \uline{On a alors les isomorphismes suivants}:
@@M15@@
\uline{où on note} @@M16@@ \uline{l'immersion de} @@M17@@ \uline{dans} @@M18@@ \uline{et} @@M19@@ \uline{l'immersion de} @@M20@@ \uline{dans} @@M21@@.

\uline{En particulier}, \uline{pour toute immersion fermée régulière} @@M22@@, \uline{on a}
@@M23@@

i) On a vu en 2.3 que @@M24@@ pouvait être muni d'une structure de @@M25@@-Algèbre; d'après 2.2, cette algèbre est alternée, car l'immersion de @@M26@@ (ou @@M27@@) dans @@M28@@ est régulière. Il existe donc un homomorphisme d'algèbres et un seul:
@@M29@@
prolongeant l'identité en degré @@M30@@.

Pour voir que c'est un isomorphisme, on peut le voir localement. On se place donc au voisinage d'un point de @@M31@@, et on considère des sections @@M32@@ de @@M33@@ dans ce voisinage, ayant les propriétés indiquées dans l'énoncé de 2.5. Posons alors @@M34@@, @@M35@@, et soient @@M36@@ l'homomorphisme défini par les @@M37@@, @@M38@@ l'homomorphisme défini par les @@M39@@, @@M40@@ et @@M41@@ les
